% ============================================================
% ch09.tex —— 第 9 章 拉格朗日的预言：子群切割术
% ============================================================
\chapter{拉格朗日的预言：子群切割术}

这一章是篇二的心脏。我们要做一件看起来人畜无害的事：把一个群当蛋糕切开。切完之后，指数的周期、素数的特权、密码锁的齿轮数，忽然全部站在了同一块基石上。这就是抽象的复利：一次证明，处处兑现。

\section{群中之群}

\begin{kaiwei}
在钟表世界 $Z_{12}$（元素 $0$ 到 $11$，加法绕圈）里，只拿 $\{0, 4, 8\}$ 这三个元素出来，只允许它们互相做加法（照旧模 $12$ 压扁）。会不会算出这三个数之外的居民？
\end{kaiwei}

动手试遍：$4+4 = 8$，$4+8 = 12\equiv 0$，$8+8 = 16\equiv 4$，$0$ 加谁都白加，$8+0 = 8$……所有结果都落在 $\{0, 4, 8\}$ 里面，一次也没越界。这三个元素\textbf{自成一国}：内部加法封闭；单位元 $0$ 在岗；每个成员的逆元在列（$4^{-1} = 8$，$8^{-1} = 4$，互为搭档；$0^{-1} = 0$）；结合律继承母国。四项体检全过——它自己就是个群，只是住在 $Z_{12}$ 的领土里。

换个候选试试：$\{0, 1, 3\}$。$1 + 3 = 4$——\textbf{越狱}，$4$ 不在集合里。封闭性阵亡，不配称国。

\begin{dingyi}[子群]
群 $(G, *)$ 的一个非空子集 $H$，若在\textbf{同一个运算} $*$ 下自己构成群（重点查封闭与逆元，单位元和结合律自动继承），就叫 $G$ 的\textbf{子群}（subgroup）。记号 $H \le G$。
\end{dingyi}

每个群都有两个"保底"子群：只有单位元的 $\{e\}$（平庸子群），和整个群自己（不正经的子群）。介于两者之间的才叫\textbf{真子群}，它们是结构的第一道指纹。顺手给"周期"上户口：$Z_{12}$ 的子群 $\{0,4,8\}$ 里有条内部规律——$4+4+4 = 12\equiv 0$，$8+8+8 = 24\equiv0$，每个成员连加三次归零。这个"$3$"是子群的大小，也是每个元素的\Keyword{归零周期}（正式名：\textbf{元素的阶}，order）。第 6 章讲"幂迟早归一"，第 8 章习题"四张牌轮换几次回家"，说的全是它。

\section{切蛋糕：盘子摆上桌}

好戏开演。拿子群 $H = \{0,4,8\}$ 当\textbf{模具}，去"扣" $Z_{12}$ 里的每个元素 $x$——所谓"扣"，就是把模具里的每个成员统统加上 $x$，扣出一个"盘子"：
\[
H + x \;:=\; \{x,\; 4+x,\; 8+x\} \pmod{12}.
\]

\begin{center}
\begin{tabular}{cl}
\toprule
用谁扣 & 扣出的盘子 \\
\midrule
$x = 0$ & $\{0, 4, 8\}$ \\
$x = 1$ & $\{1, 5, 9\}$ \\
$x = 2$ & $\{2, 6, 10\}$ \\
$x = 3$ & $\{3, 7, 11\}$ \\
$x = 4$ & $\{4, 8, 0\}$ —— 与 $x=0$ 的盘子\textbf{完全相同} \\
$x = 5, 6, \dots$ & 只是反复复现上述四个盘子 \\
\bottomrule
\end{tabular}
\end{center}

停一停，看清三件事：

\textbf{第一，盘子恰好四种}，而 $Z_{12}$ 有 $12$ 个元素。每个元素都恰好落进一个盘子（$x$ 一定在盘子 $H+x$ 里，因为 $0\in H$，$x = 0 + x$）。

\textbf{第二，每个盘子一样大}——都是 $3$ 个元素。模具多大，盘子就多大（加 $x$ 只是把模具整体平移，不增不减：$h_1 + x = h_2 + x$ 当且仅当 $h_1 = h_2$，可约去）。

\textbf{第三，盘子互不重叠}。$x=0$ 的盘子和 $x=1$ 的盘子没有共同成员，和 $x=2$、$x=3$ 的也没有——$12$ 个元素被四个盘子\textbf{不重不漏}地瓜分。

把三个数字排在一起：
\[
\underbrace{12}_{\text{群的大小}} \;=\; \underbrace{4}_{\text{盘子数}} \;\times\; \underbrace{3}_{\text{每盘元素数（}=\text{子群大小）}} .
\]
\textbf{子群的大小整除群的大小}：$3 \mid 12$。

换个群再试，看是不是巧合。三角形对称群（$6$ 元）里的"纯旋转"成员 $\{e, r, r^2\}$：封闭（转+转还是转）、有逆（$r^{-1} = r^2$）——子群，大小 $3$。拿它当模具扣全体：扣出 $H = \{e, r, r^2\}$ 和 $H' = \{s, sr, sr^2\}$（把每个翻转成员加——乘——上模具里的元素，恰好收齐三个翻转）两个盘子，$2\times3 = 6$ \checkmark，$3 \mid 6$。再试 $Z_7$ 配加法、子群 $\{0\}$：$7$ 个盘子每盘 $1$ 个，$1\mid7$ \checkmark。次次如此。这不是巧合，是定理：

\begin{dingli}[拉格朗日定理，1771]
有限群 $G$ 的任何子群 $H$，其元素个数 $|H|$ 整除 $G$ 的元素个数 $|G|$。
\end{dingli}

\begin{proof}[证明（三步，全程只用四条公理）]
\textbf{第一步：两个盘子要么全同，要么不相交。}设盘子 $H+x$ 与 $H+y$ 共享一个成员：$h_1 + x = h_2 + y$（$h_1, h_2 \in H$）。移项（群运算加法记号下移项就是两边加减元）得 $x = h_2 - h_1 + y$。注意 $h_2 - h_1 \in H$（封闭 + 逆元：$-h_1 \in H$，相加仍在）。于是 $H + x = H + (h_2 - h_1 + y)$；而"模具加自己人"原地不动（$H + h = H$，因为 $h\in H$ 时 $H+h$ 只是 $H$ 的重新排列），所以 $H + x = H + y$。\textbf{共享一个成员就共享全体}——盘子之间没有"部分重叠"这回事。

\textbf{第二步：盘子不多不少盖住全群。}每个元素 $x$ 都在盘子 $H+x$ 里（$x = 0 + x$，$0 \in H$）。全体盘子的并集就是 $G$。

\textbf{第三步：每个盘子与模具一样大。}映射 $h \mapsto h + x$ 把 $H$ 的成员一一对应地搬进盘子 $H+x$（不同 $h$ 搬进不同位置，可约去；盘子每个成员都有出处，"搬回去"就是逆映射 $z \mapsto z - x$，逆元保证 $z - x \in H$）。一一对应，所以 $|H+x| = |H|$。

三步合体：全群被若干个互不重叠、大小全为 $|H|$ 的盘子瓜分，故
\[
|G| = (\text{盘子数}) \times |H| ,
\]
右边是整数乘积，$|H|$ 整除 $|G|$。证毕。
\end{proof}

\begin{cidian}
土名"模具扣出的盘子"$\to$ \textbf{陪集}（coset）；土名"群中之群"$\to$ \textbf{子群}；土名"归零周期"$\to$ \textbf{元素的阶}（order）。盘子数有个正式名字：\textbf{指数}（index），记 $[G:H]$。拉格朗日定理比"群"这个概念早出生六十年——拉格朗日 1771 年研究代数方程时先发现了"置换的盘子现象"，六十年后伽罗瓦把"群"正式命名，定理才找到了户口。
\end{cidian}

\section{立即收割：元素的阶也整除}

拉格朗日定理的第一件赠品是推论（推论 $=$ 用定理再走半步）：

\begin{tuilun}[阶整除群阶]
有限群中任何元素 $a$ 的阶（$a$ 连乘多少次回到 $e$）整除群的大小 $|G|$。
\end{tuilun}
\begin{proof}
取元素 $a$，让它连续自我复合：$a, a^2, a^3, \dots$。全体幂 $\{a, a^2, \dots, a^{d-1}, e\}$（$d$ 为 $a$ 的阶）构成一个子群——封闭（幂乘幂还是幂，指数相加后取模 $d$ 绕回）、含 $e$、含逆（$a^k$ 的逆是 $a^{d-k}$）。它叫\textbf{循环子群}，大小恰为 $d$。拉格朗日定理：$d \mid |G|$。
\end{proof}

翻成具体语言：在 $12$ 小时的钟面上，任何时刻连加同一增量，回到原点的周期必是 $12$ 的因数（$1, 2, 3, 4, 6, 12$ 之一）——绝不会出现"周期 $5$"的时刻表。周期不是自由选的，它被群的大小"整除约束"钉死。

\section{兑现第一支欠条：费马小定理}

现在收割第 6 章的欠条。目标：$p$ 素数、$p\nmid a$ 时证 $a^{p-1}\equiv1\pmod p$。

舞台：$Z_p$ 去掉 $0$，居民 $\{1, 2, \dots, p-1\}$ 共 $p-1$ 位，运算取模 $p$ 乘法。第 7 章你已体检过它是群（逆元人人有——第 5 章的倒数名单实验早就支持了这一点，第 10 章补发正式证明）。

任取居民 $a$，看它的阶 $d$（$a^d \equiv 1$ 的最小正指数）。由上一节推论：
\[
d \mid (p-1), \qquad\text{即}\qquad p - 1 = d\times k \ \text{（}k\text{ 为整数）}.
\]
于是
\[
a^{p-1} = a^{dk} = \left(a^{d}\right)^{k} \equiv 1^{k} = 1 \pmod p .
\]
\textbf{证毕。}三行。

回头看第 1 章那张个位数字周期表，现在全懂了：底数 $2$ 在模 $10$ 世界里的周期 $4$，其实是"$2$ 在群里的阶"；表中底数的周期只有 $1, 2, 4$（都是 $4$ 的因数），因为它们活动的"单位群"大小是 $4$（$Z_{10}$ 的居民里只有 $1,3,7,9$ 有倒数，凑成 $4$ 元群——第 10 章正式介绍）。\textbf{周期必整除群阶}这一句拉格朗日语，翻译成人话就是那张表的全部规律。

再看第 6 章的"齿轮数"$(p-1)(q-1)$：RSA 的加密指数 $e$ 和解密指数 $d$ 满足 $ed \equiv 1 \pmod{(p-1)(q-1)}$——为什么齿轮偏偏是 $(p-1)(q-1)$？因为"与 $N$ 互素的剩余类"构成大小恰为 $(p-1)(q-1)$ 的乘法群（$p-1$ 位居民来自模 $p$ 世界、$q-1$ 位来自模 $q$ 世界，中国剩余定理把它们编织成一个 $(p-1)(q-1)$ 人的大群——数论里管这个大小叫\textbf{欧拉函数} $\varphi(N)$）。幂在这个群里转圈，圈长整除 $\varphi(N)$，所以"先转 $e$ 步再转 $d$ 步"共 $ed$ 步 $\equiv 1 \pmod{\varphi}$ 步，回到原点——\textbf{加密再解密必还原，是拉格朗日定理在给密码锁上保险}。（这套拼装的最后一颗螺丝——"转 $1$ 圈回到原点 $\Rightarrow$ 任意元素转整圈回到自己"——在第 10 章拧紧。）

\section{拉格朗日的预言}

拉格朗日本人 1771 年做切割研究时，真正的猎物是\textbf{代数方程}。背景：二次方程有求根公式（初中学过），三次、四次也有（十六世纪意大利数学家卡丹、费拉里在吵架、决斗与剽窃指控中搞出来的——数学史上最狗血的一章）；五次呢？拉格朗日分析根的置换时隐约看到了天花板，他在论文里写下了那句著名的预感："这件事看来是在向数学家的机智挑战。"

答案半个世纪后揭晓：1824 年，$22$ 岁的挪威青年阿贝尔证明\textbf{五次及以上的一般方程不存在根根公式}；1832 年，$20$ 岁的法国青年伽罗瓦更进一步——他给\textbf{每个具体方程}配了一个群（根之间的对称群），证明"有没有求根公式"取决于这个群能不能被切成特定形状的塔（\textbf{可解群}）。拉格朗日的盘子、本章的切割术，在伽罗瓦手里变成了成套的解剖刀。完整故事（含两位天才的悲剧结局）在第 12 章选读。

\begin{dongshou}
\begin{itemize}
  \yisi 在三角形对称群（$6$ 元）里找出全部子群，逐个验证"阶整除 $6$"。（答案骨架：$1$ 元 $\{e\}$；$2$ 元 $\{e,s\}, \{e,sr\}, \{e,sr^2\}$——三个；$3$ 元 $\{e,r,r^2\}$；$6$ 元全群。$6$ 的因数恰是 $1,2,3,6$，一个不多一个不少。特别体会：\textbf{$4$ 元和 $5$ 元子群不可能存在}——$4\nmid6$，$5\nmid 6$，拉格朗日直接否决，连找都不用找。）
  \yier 用"阶整除群阶"秒答：$Z_{13}$ 去零的乘法群有 $12$ 位居民，任何元素的阶只能是 $1, 2, 3, 4, 6, 12$。由此求 $2^{100} \bmod 13$：先试出 $2$ 的阶（$2^2=4$，$2^3=8$，$2^4=16\equiv3$，$2^6\equiv 64\equiv 12 \equiv -1$，所以 $2^{12}\equiv1$ 且没有更小的 $d$ 让 $2^d \equiv 1$？——注意 $2^6\equiv -1$ 排除了 $d\mid 6$，还剩 $d = 4$ 已试 $\equiv3$，故阶恰为 $12$）；于是 $100 = 12\times8 + 4$，$2^{100}\equiv 2^4 \equiv 3 \pmod{13}$。
  \yier 一个 $5$ 元群里可能有什么阶的元素？（阶只能整除 $5$：即 $1$ 或 $5$。除单位元外人人阶 $5$——这样的群叫"素数阶循环群"，下一章它是"域"的最佳演员。）顺手证明：素数阶群一定是循环群（任取非单位元 $a$，它生成的循环子群阶整除素数 $p$ 且不是 $1$，只能等于 $p$——子群挤满全群）。
\end{itemize}
\end{dongshou}

\begin{shaonao}
扑克魔术师的最爱——\textbf{完美洗牌}（faro shuffle）：$52$ 张牌切两半 $26+26$，一张隔一张精确交错插回。这种洗法藏着一个秘密：\textbf{重复若干次后牌序完全复原}。跟踪每张牌的新位置（顶牌记位置 $0$）：一次完美洗牌后，位置 $j$ 的牌跳到位置 $2j \bmod 51$（模 $51$，因为除了最后一张"钉子户"，全体在绕一个 $51$ 位置的圈）。于是洗 $k$ 次后，牌来到位置 $2^k j \bmod 51$——全体复原当且仅当 $2^k \equiv 1 \pmod{51}$。用本章工具求最小的 $k$：$2^8 = 256 = 5\times51+1$，\textbf{$8$ 次外洗完全复原}！（魔术师据此设计"洗八次牌序自动归位"的招牌戏法。）顺带用拉格朗日验算：$51 = 3\times17$，与 $51$ 互素的居民有 $\varphi(51) = 2\times16 = 32$ 位，$2$ 的阶 $8$ 整除 $32$ \checkmark。你在牌桌上见到了本章定理的现场演出。
\end{shaonao}

\begin{chengshi}
第一步证明里"模具加自己人原地不动"（$H + h = H$）依赖 $H$ 的封闭与逆元（$h_1 + h \in H$ 且 $(h_1+h) - h = h_1 \in H$），我们把验证嵌在了行文中；乘法记号的群（如对称群）请自行翻译（$Hx$、$h_1h_2^{-1}\in H$ 等）。另一个重要提醒：拉格朗日定理的\textbf{逆命题不成立}——"$d$ 整除 $|G|$"不保证存在 $d$ 阶子群，最小反例是 $12$ 元的交错群 $A_4$（$6\mid 12$ 但 $A_4$ 没有 $6$ 元子群）；不过对"循环群"逆命题成立（每个因数都有对应子群），这也是数论里 $\varphi$ 函数好用的原因。深水区留给大学。
\end{chengshi}

蛋糕切完，工具栏新增：子群、陪集、元素的阶、拉格朗日定理。数论的两笔大欠条（费马小定理、周期规律）已经全部还清。下一章，我们处理一个悬了三章的经济问题：\textbf{在什么世界里，除法才全面通行？}——以及亲眼围观"分数的诞生"。
